\subsection{准素分解的存在性}

定义 2. 设 A 是诺特环，M 是 A-模，N 是 M 的子模。M 的一族有限个关于 M 为准素的子模 \((\mathbf{Q}_i)_{i\in I}\)，若满足 \(\mathbf{N} = \bigcap_{i\in I}\mathbf{Q}_i\)，则称为 N 在 M 中的准素分解。

例。取 \(\mathbf{A} = \mathbf{Z}\)，\(\mathbf{M} = \mathbf{Z}\)，\(\mathbf{N} = n\mathbf{Z}\)，其中整数 \(n > 0\)。若 \(n = p_{1}^{\alpha_1} \ldots p_{k}^{\alpha_k}\) 是 \(n\) 的素因子分解，

\[
n \mathbf {Z} = (p _ {1} ^ {\alpha_ {1}} \mathbf {Z}) \cap \dots \cap (p _ {k} ^ {\alpha_ {k}} \mathbf {Z})
\]

这是 \(n\mathbf{Z}\) 在 \(\mathbf{Z}\) 中的准素分解，由第 1 节例 4 得出。

滥用语言地说，关系\(\mathbf{N} = \bigcap_{i \in I} \mathbf{Q}_i\)称为\(\mathbf{N}\)中的准素分解\(\mathbf{M}\)。同样也可说，\(\{0\} = \bigcap_{i \in I} (\mathbf{Q}_i / \mathbf{N})\)是\(\{0\}\)中的准素分解\(\mathbf{M} / \mathbf{N}\)。如果\((\mathbf{Q}_i)_{i \in I}\)是\(\mathbf{N}\)中的准素分解\(\mathbf{M}\)，典范映射从\(\mathbf{M} / \mathbf{N}\)到\(\bigoplus_{i \in I} (\mathbf{M} / \mathbf{Q}_i)\)是单射。反之，设\(\mathbf{N}\)是\(\mathbf{M}\)的子模，且\(f\)是从\(\mathbf{M} / \mathbf{N}\)到有限直和\(\mathbf{P} = \bigoplus_{i \in I} \mathbf{P}_i\)，其中每个集合\(\operatorname{Ass}(\mathbf{P}_i)\)都归结为单个元素\(p_i\)；令\(f_i\)是同态\(\mathbf{M} / \mathbf{N} \to \mathbf{P}_i\)，它由\(f\)与投影\(\mathbf{P} \to \mathbf{P}_i\)的复合得到，并令\(Q_i / \mathbf{N}\)是\(f_i\)的核；则与\(Q_i\)不同的\(\mathbf{M}\)关于\(\mathbf{M}\)都是准素的（第 1 节，定义 1），且\(\mathbf{N} = \bigcap_{i \in I} Q_i\)。此外，\(\operatorname{Ass}(\mathbf{M} / \mathbf{N}) \subset \bigcup_{i \in I} \{\mathfrak{p}_i\}\)（根据第 1 节，第 1 节，命题 3 的推论 2）。

定理 1. 设 M 是诺特环上的有限生成模，N 是 M 的子模。存在如下形式的 N 的准素分解

\[
N = \bigcap_ {\mathfrak {p} \in \operatorname{Ass} (M / N)} Q (\mathfrak {p})\tag{1}
\]

其中，对于所有 \(\mathfrak{p} \in \operatorname{Ass}(\mathbf{M} / \mathbf{N})\)，\(\mathbf{Q}(\mathfrak{p})\) 是 \(\mathfrak{p}\)-准素的（相对于 \(\mathbf{M}\)）。

我们可以用 M/N 替换 M，因而设 N = 0。根据 § 1，第 4 节，定理 2 的推论，Ass(M) 是有限集；根据 § 1，第 1 节，命题 4，对于每个 \(p \in \text{Ass}(M)\)，存在 M 的一个子模 \(Q(p)\)，使得 \(\text{Ass}(M/Q(p)) = \{p\}\) 且 \(\text{Ass}(Q(p)) = \text{Ass}(M) - \{p\}\)。记 \(P = \bigcap_{p \in \text{Ass}(M)} Q(p)\)；对于所有 \(p \in \text{Ass}(M)\)，有 \(\text{Ass}(P) \subset \text{Ass}(Q(p))\)，从而 \(\text{Ass}(P) = \varnothing\)，这意味着 P = 0（§ 1，第 1 节，命题 2 的推论 1），因而定理得证。

